\section{Část II: Společná informatika}

{\small Společná informatika je floor pro společné okruhy: cíl je umět každé téma
bez nuly a rychle poskládat zelenou odpověď. Tady stačí kostra, definice, jeden příklad
a hlavní pasti.}

% ============================================================
\subsection{Automaty a jazyky}
\priorita{střední}{zelená (2 body)}

\begin{tema}
\textbf{Automaty a jazyky}
\begin{itemize}[leftmargin=1.2em,itemsep=0pt,topsep=1pt]
  \item Regulární jazyky
  \item regulární gramatiky
  \item deterministický a nedeterministický konečný automat
  \item regulární výrazy
  \item Bezkontextové jazyky
  \item bezkontextové gramatiky, jazyk generovaný gramatikou
  \item zásobníkový automat, třída jazyků přijímaných zásobníkovými automaty
  \item Rekurzivně spočetné jazyky
  \item gramatika typu 0
  \item Turingův stroj
  \item algoritmicky nerozhodnutelné problémy
  \item Chomského hierarchie
  \item schopnost zařazení konkrétního jazyka do Chomského hierarchie (zpravidla sestrojení odpovídajícího automatu či gramatiky)
\end{itemize}
\end{tema}

\ucivo
\begin{intuice}
Jazyk je množina slov. Čím silnější model výpočtu dovolím, tím složitější jazyky umím
popsat: regulární přes konečnou paměť, bezkontextové přes zásobník, rekurzivně spočetné
přes Turingův stroj.
\end{intuice}

\begin{center}
\begin{minipage}{0.43\linewidth}
\centering
\begin{tikzpicture}[>=Stealth, node distance=12mm,
  every node/.style={circle,draw,minimum size=7mm,font=\scriptsize,inner sep=0pt}]
  \node (q0) {$q_0$};
  \node[double,right=of q0] (q1) {$q_1$};
  \draw[->] (-0.9,0) -- (q0);
  \draw[->] (q0) edge[loop above] node[draw=none,font=\scriptsize]{$b$} (q0);
  \draw[->] (q0) -- node[above,draw=none,font=\scriptsize]{$a$} (q1);
  \draw[->] (q1) edge[loop above] node[draw=none,font=\scriptsize]{$a,b$} (q1);
\end{tikzpicture}\\[-1pt]
{\scriptsize DFA pro slova obsahující aspoň jedno $a$.}
\end{minipage}\hfill
\begin{minipage}{0.49\linewidth}
\centering
\begin{tikzpicture}[scale=0.82, every node/.style={font=\scriptsize}]
  \draw[fill=defblue!6,draw=defblue] (0,0) ellipse (2.45cm and 1.45cm);
  \draw[fill=intugreen!10,draw=intugreen] (0,0) ellipse (1.85cm and 1.05cm);
  \draw[fill=orange!12,draw=orange!80!black] (0,0) ellipse (1.18cm and 0.68cm);
  \draw[fill=gapred!8,draw=gapred] (0,0) ellipse (0.55cm and 0.32cm);
  \node at (0,0) {reg.};
  \node at (0,0.63) {bezkontextové};
  \node at (0,1.06) {kontextové};
  \node at (0,1.43) {typ 0 / RE};
\end{tikzpicture}\\[-1pt]
{\scriptsize Chomského hierarchie jako vnořené třídy.}
\end{minipage}
\end{center}

\begin{odpoved}
\setlength{\parskip}{2pt}
\noindent\bod{1}\ \textbf{Minimum:} \emph{DFA} je pětice $(Q,\Sigma,\delta,q_0,F)$ s totální funkcí $\delta:Q\times\Sigma\to Q$; \emph{NFA} dovolí množinu přechodů a $\varepsilon$-kroky.
\emph{Regulární jazyk} je jazyk přijímaný konečným automatem, ekvivalentně popsatelný regulárním výrazem nebo regulární gramatikou.
\emph{Turingův stroj} má konečné řízení, pásku a přechodovou funkci; přijímá vstupy, na nichž zastaví v přijímajícím stavu.\par
\noindent\bod{2}\ \textbf{+ k tomu:} \emph{Bezkontextový jazyk} generuje CFG s pravidly
$A\to\alpha$ a přijímá ho \emph{PDA}, tedy konečný automat se zásobníkem. \emph{Rekurzivně spočetný jazyk} přijímá Turingův stroj.
Zásobník dá paměť pro vnoření, například
$\{a^n b^n\mid n\ge0\}$, ale ne pro dvě nezávislá počítání jako $a^n b^n c^n$.\par
\noindent\bod{3}\ \textbf{Bonus:} Typ 0 gramatiky odpovídají rekurzivně spočetným jazykům
a Turingovým strojům. Existují algoritmicky nerozhodnutelné problémy, klasicky problém
zastavení. U zařazování jazyka zkus sestrojit automat nebo gramatiku, případně říct, proč
nižší třída nestačí. Pumping lemma a Myhill-Nerodeova věta slouží k důkazu neregularity;
pro CFL existuje pumping lemma a Chomského normální tvar. PDA přijímá koncovým stavem
nebo prázdným zásobníkem, obě varianty jsou ekvivalentní. Regulární jazyky jsou uzavřené
na doplněk, průnik a sjednocení, CFL ne na průnik obecně.
\end{odpoved}

\begin{recall}
\begin{itemize}[leftmargin=1.2em,itemsep=0pt,topsep=1pt]
  \item Čím se liší DFA a NFA a proč mají stejnou vyjadřovací sílu?
  \item Jaký automat přijímá bezkontextové jazyky a k čemu je zásobník?
  \item Kam v hierarchii patří regulární, bezkontextové a rekurzivně spočetné jazyky?
  \item Uveď příklad nerozhodnutelného problému.
\end{itemize}
\end{recall}
\past{Neříkej, že NFA je silnější než DFA. Je pohodlnější, ale pro konečné automaty popisuje stejnou třídu jazyků.}

% ============================================================
\subsection{Algoritmy a datové struktury}
\priorita{střední}{zelená (2 body)}

\begin{tema}
\textbf{Algoritmy a datové struktury}
\begin{itemize}[leftmargin=1.2em,itemsep=0pt,topsep=1pt]
  \item Časová složitost algoritmů
  \item časová a prostorová složitost algoritmu
  \item měření velikosti dat
  \item složitost v nejlepším, nejhorším a průměrném případě
  \item asymptotická notace
  \item Třídy složitosti
  \item třídy P a NP
  \item převoditelnost problémů, NP-těžkost a NP-úplnost
  \item příklady NP-úplných problémů a převodů mezi nimi
  \item Metoda rozděl a panuj
  \item princip rekurzivního dělení problému na podproblémy
  \item výpočet složitosti pomocí rekurentních rovnic
  \item Master theorem (kuchařková věta) (bez důkazu)
  \item aplikace
  \item Mergesort
  \item násobení dlouhých čísel
  \item Binarní vyhledávací stromy
  \item definice vyhledávacího stromu
  \item operace s nevyvažovanými stromy
  \item AVL stromy (definice)
  \item Třídění
  \item primitivní třídicí algoritmy (Bubblesort, Insertsort)
  \item Quicksort
  \item dolní odhad složitosti porovnávacích třídicích algoritmů
  \item Grafové algoritmy
  \item prohledávání do šířky a do hloubky
  \item topologické třídění orientovaných grafů
  \item nejkratší cesty v ohodnocených grafech (Dijkstrův a Bellmanův-Fordův algoritmus)
  \item minimální kostra grafu (Jarníkův a Borůvkův algoritmus)
  \item toky v sítích (Ford-Fulkerson algoritmus)
\end{itemize}
\end{tema}

\ucivo
\begin{intuice}
U algoritmů se zkouší hlavně rozpoznat správnou techniku, říct invariant nebo datovou
strukturu a umět odhadnout složitost. U společné části stačí mít mapu pojmů a typické
časy v ruce.
\end{intuice}

\begin{center}
\begin{tikzpicture}[level distance=9mm,every node/.style={circle,draw,minimum size=6mm,font=\scriptsize,inner sep=0pt},
  level 1/.style={sibling distance=24mm},level 2/.style={sibling distance=12mm}]
\node {8}
  child {node {4}
    child {node {2}}
    child {node {6}}}
  child {node {12}
    child {node {10}}
    child {node {14}}};
\node[draw=none,anchor=west,font=\scriptsize,align=left] at (3.1,-0.8)
  {BST: vlevo menší, vpravo větší\\AVL: v každém uzlu rozdíl výšek nejvýše 1};
\end{tikzpicture}
\end{center}

\begin{odpoved}
\setlength{\parskip}{2pt}
\noindent\bod{1}\ \textbf{Minimum:} $f\in O(g)$ znamená $\exists c,n_0\ \forall n\ge n_0:\ f(n)\le c g(n)$; $\Omega$ je dolní a $\Theta$ oboustranný asymptotický odhad.
\emph{P} je třída rozhodovacích problémů řešitelných v polynomiálním čase deterministickým strojem.
\emph{NP} je třída problémů, jejichž ano-certifikát lze ověřit v polynomiálním čase.\par
\noindent\bod{2}\ \textbf{+ k tomu:} \emph{Polynomiální redukce} $A\le_p B$ je polynomiálně spočitatelná transformace $x\mapsto f(x)$ se $x\in A \iff f(x)\in B$.
\emph{NP-těžký} problém je takový, že se na něj polynomiálně redukuje každý problém z NP; \emph{NP-úplný} je NP-těžký a leží v NP.
\emph{BST} má v každém uzlu vlevo menší a vpravo větší klíče; \emph{AVL} je BST, kde se výšky levého a pravého podstromu v každém uzlu liší nejvýše o 1.
Rozděl a panuj: rozděl,
vyřeš podproblémy, slož výsledek, například mergesort $T(n)=2T(n/2)+O(n)=O(n\log n)$.
AVL drží vyvážení rotacemi. Porovnávací třídění má dolní
odhad $\Omega(n\log n)$.\par
\noindent\bod{3}\ \textbf{Bonus:} Quicksort má průměrně $O(n\log n)$ a nejhůře $O(n^2)$.
BFS je fronta a vzdálenosti v neohodnoceném grafu, DFS je zásobník a časy průchodu.
Topologické třídění funguje pro DAG. Dijkstra vyžaduje nezáporné hrany, Bellman-Ford
zvládá záporné hrany, Jarník a Borůvka řeší MST, Ford-Fulkerson zlepšuje tok po
reziduálních cestách. Master theorem pro $T(n)=aT(n/b)+\Theta(n^c)$ porovnává
$a/b^c$; Karacuba násobí dlouhá čísla v $O(n^{\log_2 3})$. NP převody jsou
reflexivní a tranzitivní; typické NP-úplné problémy jsou SAT, 3-SAT, klika,
nezávislá množina, $k$-obarvitelnost, Hamiltonovská kružnice, batoh a součet podmnožiny.
AVL má hloubku $\Theta(\log n)$, porovnávací třídění má dolní mez $\Omega(n\log n)$.
\end{odpoved}

\begin{recall}
\begin{itemize}[leftmargin=1.2em,itemsep=0pt,topsep=1pt]
  \item Co musí platit, aby byl problém NP-úplný?
  \item Jak spočítáš mergesort přes rekurentní rovnici?
  \item Jaký je rozdíl mezi BST a AVL stromem?
  \item Kdy použiješ Dijkstru a kdy Bellman-Forda?
\end{itemize}
\end{recall}
\past{NP-těžký problém nemusí být v NP. NP-úplnost vyžaduje obojí: patří do NP a každý problém z NP se na něj polynomiálně převede.}

% ============================================================
\subsection{Programovací jazyky}
\priorita{střední}{zelená (2 body)}

\begin{tema}
\textbf{Programovací jazyky}

Některé následující body definují varianty požadavků pro různé individuální volby povinně volitelných předmětů.
Vyžaduje se zvládnutí všech bodů bez označení \textcircled{V} a zvládnutí všech bodů s označením \textcircled{V} pro jeden z jazyků C\#, C++ nebo Java.
\begin{itemize}[leftmargin=1.2em,itemsep=0pt,topsep=1pt]
  \item Koncepty pro abstrakci, zapouzdření a polymorfismus.
  \item související konstrukty programovacích jazyků
  \item třídy, rozhraní, metody, datové položky, dědičnost, viditelnost
  \item (dynamický) polymorfismus, statické a dynamické typování
  \item jednoduchá dědičnost
  \item \textcircled{V} virtuální a nevirtuální metody v C++ a C\#
  \item vícenásobná dědičnost a její problémy
  \item \textcircled{V} interfaces v C\#
  \item implementace rozhraní (interface)
  \item vhodné použití uvedených konceptů
  \item Primitivní a objektové typy a jejich reprezentace.
  \item číselné a výčtové typy
  \item \textcircled{V} hodnotové a referenční typy v C\#
  \item \textcircled{V} reference, imutabilní typy a boxing v C\# a Javě
  \item Generické typy a funkcionální prvky (procedurálních programovacích jazyků).
  \item \textcircled{V} generické typy v Javě a C\# (bez omezení typových parametrů)
  \item \textcircled{V} typy reprezentující funkce v C++, C\#, nebo Javě
  \item lambda funkce a funkcionální rozhraní
  \item Manipulace se zdroji a mechanizmy pro ošetření chyb.
  \item správa životního cyklu zdrojů v případě výskytu chyb
  \item \textcircled{V} using v C\#
  \item konstrukce pro obsluhu a propagaci výjimek
  \item Životní cyklus objektů a správa paměti.
  \item alokace (alokace statická, na zásobníku, na haldě)
  \item inicializace (konstruktory, volání zděděných konstruktorů)
  \item destrukce (destruktory, finalizátory)
  \item explicitní uvolňování objektů, reference counting, garbage collector
  \item Vlákna a podpora synchronizace.
  \item reprezentace vláken v programovacích jazycích
  \item specifikace funkce vykonávané vláknem a základní operace na vlákny
  \item časově závislé chyby a mechanizmy pro synchronizaci vláken
  \item Implementace základních prvků objektových jazyků.
  \item základní objektové koncepty v konkrétním jazyce
  \item implementace a interní reprezentace primitivních typů
  \item implementace a interní reprezentace složených typů a objektů
  \item implementace dynamického polymorfismu (tabulka virtuálních metod)
  \item Nativní a interpretovaný běh, řízení překladu a sestavení programu.
  \item reprezentace programu, bytecode, interpret jazyka
  \item just-in-time (JIT) a ahead-of-time (AOT) překlad
  \item proces sestavení programu, oddělený překlad, linkování
  \item staticky a dynamicky linkované knihovny
  \item běhové prostředí procesu a vazba na operační systém
\end{itemize}
\end{tema}

\ucivo
\begin{intuice}
Drž se C\#. Odpověď stav na třech vrstvách: jazykové konstrukty, běhová reprezentace
a bezpečné zacházení se zdroji a chybami.
\end{intuice}

\begin{odpoved}
\setlength{\parskip}{2pt}
\noindent\bod{1}\ \textbf{Minimum:} Abstrakce skrývá detail, zapouzdření chrání stav
přes viditelnost, polymorfismus dovolí volat stejnou operaci nad různými typy. V C\#
jsou třídy, rozhraní, metody, položky, jednoduchá dědičnost tříd a vícenásobná
implementace rozhraní. Typování je statické, dynamická vazba se používá u virtuálních
metod.\par
\noindent\bod{2}\ \textbf{+ k tomu:} Hodnotové typy se typicky ukládají přímo a kopírují
hodnotu, referenční typy se obsluhují přes referenci na objekt na haldě. Boxing zabalí
hodnotový typ do objektu. Generika v C\# parametrizují třídy a metody typem, delegáti
a lambda výrazy reprezentují funkce. Výjimky se propagují zásobníkem, \texttt{using}
zajistí \texttt{Dispose} i při chybě.\par
\noindent\bod{3}\ \textbf{Bonus:} Konstruktor inicializuje objekt a volá konstruktor
předka, správu paměti řeší garbage collector, finalizátor není náhrada za uvolnění
externího zdroje. Dynamický polymorfismus bývá přes tabulku virtuálních metod. C\#
běží nad CLR, program se překládá do IL, typicky JIT do strojového kódu, knihovny
mohou být staticky nebo dynamicky navázané. Vlákna sdílejí paměť, proto hrozí race
conditions a používají se zámky. Diamantový problém řeší C\# vícenásobnou implementací
rozhraní místo vícenásobné dědičnosti tříd. \texttt{ref}, \texttt{out} a \texttt{in}
předávají argument odkazem (u \texttt{in} jen pro čtení); fungují pro hodnotové i
referenční typy. \texttt{readonly} a schovaný setter pomáhají s imutabilitou.
RAII v C++ váže zdroj na život objektu; \texttt{using} je \texttt{try/finally}
s \texttt{Dispose}.
\end{odpoved}

\begin{recall}
\begin{itemize}[leftmargin=1.2em,itemsep=0pt,topsep=1pt]
  \item Jaký je rozdíl mezi hodnotovým a referenčním typem v C\#?
  \item Co v C\# znamená \texttt{virtual}, \texttt{override} a interface?
  \item Kdy použiješ \texttt{using} a proč nestačí garbage collector?
  \item Jak spolu souvisí IL, CLR a JIT?
\end{itemize}
\end{recall}
\past{V C\# není vícenásobná dědičnost tříd. Vícenásobně lze implementovat rozhraní.}

% ============================================================
\subsection{Architektura počítačů a operačních systémů}
\priorita{střední}{zelená (2 body)}

\begin{tema}
\textbf{Architektura počítačů a operačních systémů}
\begin{itemize}[leftmargin=1.2em,itemsep=0pt,topsep=1pt]
  \item Základní architektura počítače, reprezentace čísel, dat a programů.
  \item reprezentace a přístup k datům v paměti, adresa, adresový prostor
  \item ukládání jednoduchých a složených datových typů
  \item základní aritmetické a logické operace
  \item Instrukční sada, vazba na prvky vyšších programovacích jazyků.
  \item Implementovat běžné programové konstrukce vyšších jazyků (přiřazení, podmínka, cyklus, volání funkce) pomocí instrukcí procesoru
  \item Zapsat běžnou konstrukci vyššího jazyka (přiřazení, podmínka, cyklus, volání funkce), která odpovídá zadané sekvenci (vysvětlených) instrukcí procesoru
  \item Podpora pro běh operačního systému.
  \item privilegovaný a neprivilegovaný režim procesoru
  \item jádro operačního systému
  \item Rozhraní periferních zařízení a jejich obsluha.
  \item Popsat roli řadiče zařízení při programem řízené obsluze zařízení (PIO), pro zadané adresy a funkce vstupních a výstupních portů implementovat programem řízenou obsluhu zadaného zařízení (myš, disk)
  \item Popsat roli přerušení při programem řízené obsluze zařízení (PIO), na úrovni vykonávání instrukcí popsat reakci procesoru (hardware) a operačního systému (software) na žádost o přerušení
  \item Základní abstrakce, rozhraní a mechanizmy OS pro běh programů, sdílení prostředků a vstup/výstup.
  \item neprivilegované (uživatelské) procesy
  \item sdílení procesoru
  \item procesy, vlákna, kontext procesu a vlákna
  \item přepínání kontextu, kooperativní a preemptivní multitasking
  \item plánování běhu procesů a vláken, stavy vlákna
  \item sdílení paměti
  \item Vysvětlit rozdíl mezi virtuální a fyzickou adresou a identifikovat, zda se v zadaném kontextu či fragmentu kódu používá virtuální nebo fyzická adresa
  \item Na zadaném příkladu identifikovat a vysvětlit význam komponent virtuální a fyzické adresy (číslo stránky, číslo rámce, offset)
  \item Pro konkrétní adresy a obsah jednoúrovňové stránkovací tabulky řešit úlohy překladu adres
  \item Vysvětlit roli virtuálních adresových prostorů v ochraně paměti procesů a vláken
  \item sdílení úložného prostoru
  \item soubory, analogie s adresovým prostorem
  \item abstrakce a rozhraní pro práci se soubory
  \item Paralelismus, vlákna a rozhraní pro jejich správu, synchronizace vláken.
  \item časově závislé chyby (race conditions)
  \item kritická sekce, vzájemné vyloučení
  \item základní sychronizační primitiva, jejich rozhraní a použití
  \item zámky
  \item aktivní a pasivní čekání
\end{itemize}
\end{tema}

\ucivo
\begin{intuice}
Procesor vykonává instrukce nad registry a pamětí, OS odděluje uživatelský kód od jádra
a dává procesům iluze: vlastní CPU, vlastní adresový prostor a soubory jako pojmenované
úložiště.
\end{intuice}

\begin{center}
\begin{tikzpicture}[>=Stealth, font=\scriptsize]
  \node[draw,minimum width=18mm,minimum height=7mm] (va) at (0,0) {virtuální adresa};
  \node[draw,minimum width=12mm,minimum height=7mm,right=5mm of va] (pn) {stránka};
  \node[draw,minimum width=12mm,minimum height=7mm,right=0mm of pn] (off) {offset};
  \node[draw,minimum width=18mm,minimum height=7mm,right=9mm of off] (pt) {tabulka};
  \node[draw,minimum width=12mm,minimum height=7mm,right=9mm of pt] (fr) {rámec};
  \node[draw,minimum width=12mm,minimum height=7mm,right=0mm of fr] (off2) {offset};
  \node[draw,minimum width=18mm,minimum height=7mm,right=5mm of off2] (pa) {fyzická adresa};
  \draw[->] (va) -- (pn);
  \draw[->] (off) -- (pt);
  \draw[->] (pn) -- (pt);
  \draw[->] (pt) -- node[above,draw=none] {rámec} (fr);
  \draw[->] (off2) -- (pa);
  \draw[->] (fr) -- (pa);
  \draw[->] (off.south) to[out=-70,in=-110] node[below,draw=none] {beze změny} (off2.south);
\end{tikzpicture}
\end{center}

\begin{odpoved}
\setlength{\parskip}{2pt}
\noindent\bod{1}\ \textbf{Minimum:} Počítač má CPU, registry, paměť a zařízení. Data i program
jsou v paměti jako bity, adresa určuje místo v adresovém prostoru. Instrukce realizují
přiřazení, větvení, cyklus a volání funkce přes čtení, zápis, skoky a práci se zásobníkem.
OS běží v privilegovaném režimu, procesy běží uživatelsky.\par
\noindent\bod{2}\ \textbf{+ k tomu:} Řadič zařízení vystavuje porty nebo registry. PIO znamená,
že procesor programově čte a zapisuje stav a data zařízení. Přerušení dovolí zařízení
vyžádat obsluhu: hardware uloží minimální stav, skočí na obsluhu v jádře, OS uloží zbytek
kontextu a po obsluze se pokračuje.
Proces má vlastní adresový prostor a zdroje, vlákno je běžící tok s vlastním kontextem.
Preemptivní multitasking přepíná OS, kooperativní čeká na dobrovolné předání.\par
\noindent\bod{3}\ \textbf{Bonus:} Virtuální adresa se dělí na číslo stránky a offset,
tabulka stránek mapuje stránku na rámec, offset se přenese. Virtuální adresové prostory
chrání procesy mezi sebou. Soubory jsou pojmenované úložiště s rozhraním open, read,
write, close. Race condition vzniká při závislosti na proložení vláken, kritickou sekci
chrání vzájemné vyloučení, typicky zámek. Aktivní čekání točí CPU, pasivní čekání vlákno
uspí. Harvardská architektura odděluje paměť kódu a dat, von Neumann má jednu paměť.
Čísla se ukládají binárně, záporná typicky dvojkovým doplňkem, float jako znaménko,
mantisa a exponent; endianita určuje pořadí bytů a zarovnání vkládá mezery do struktur.
Syscall je řízený přechod do jádra; DMA přenáší data zařízení bez kopírování CPU.
MMU překládá stránky a page fault řeší chybějící mapování. Scheduler používá stavy
created, ready, running, blocked, terminated; algoritmy zahrnují FCFS, round robin,
multilevel feedback queue a CFS. Semafor má čítač, mutex je binární zámek, barrier
čeká na více vláken.
\end{odpoved}

\begin{recall}
\begin{itemize}[leftmargin=1.2em,itemsep=0pt,topsep=1pt]
  \item Jak instrukcemi vyjádříš podmínku, cyklus a volání funkce?
  \item Co se stane při přerušení z pohledu procesoru a OS?
  \item Jak přeložíš virtuální adresu přes jednoúrovňovou stránkovací tabulku?
  \item Jaký je rozdíl mezi procesem, vláknem a jejich kontextem?
  \item Čím se liší aktivní a pasivní čekání?
\end{itemize}
\end{recall}
\past{Virtuální adresa není fyzické místo v RAM. Je to adresa v prostoru procesu, kterou MMU a tabulka stránek překládají na fyzický rámec.}
